\newpage % The Meta-Review should at least start on a new column
\glsresetall

% Use \appendices and not \appendix due to IEEEtran.cls quirks
\appendices % if not used earlier


\nobalance
\section{Meta-Review}

The following meta-review was prepared by the program committee for the 2025 IEEE Symposium on Security and Privacy (S\&P) as part of the review process as detailed in the call for papers.

\subsection{Summary}

This paper presents \thename, a system that obfuscates (at specific learned intervals) the memory of a \gls{cvm}. The system leverages \gls{oram}. The challenge is leveraging \gls{oram} in a practical fashion to entire VM executions. To solve these, the paper leverages a unikernel design (based on Unikraft) and an adaptive exit-rate-based algorithm. The system incurs 2-4 times overhead during execution of real-world programs.



\subsection{Scientific Contributions}
\begin{itemize}
\item Provides a valuable step forward in an established field
\item Creates a new tool to enable future science
\item Establishes a new research direction
\end{itemize}

\subsection{Reasons for Acceptance}
\begin{enumerate}
\item Side-channels remain a problem for \glspl{cvm}. The reviewers appreciated how the paper positioned itself in the domain of side-channel mitigations for existing TEEs like SGX, especially ones that have leveraged \gls{oram}. The system designed by \thename helps advance our understanding of \gls{oram}'s application into whole VM TEEs.
\item The system designed in this paper can help other researchers develop ORAM-specific mitigations for \glspl{cvm}.
\end{enumerate}

\subsection{Noteworthy Concerns} % Exclude if your meta-review does not have noteworthy concerns
%\begin{enumerate} % Enumerate environment is not necessary if there is only one
The paper leaves core features (including multi-threading support and defense against other side-channels like cache) as part of future work or out-of-scope.
% \end{enumerate}
